\lesson{4}{Jan 10 2022 Mon (16:55:38)}{Civil Rights}{Unit 4}

\subsubsection{Protecting Civil Rights}

The Declaration of Independence and the U.S. Constitution provide a strong
foundation for the protection of U.S. citizen's rights. People have used
language from each to support arguments that certain laws or practices unfairly
restrict their rights.

Early Americans emphasized the last paragraph in the Declaration of
Independence, the one proclaiming that split from Britain. Over time, Americans
have focused more on the idea of rights that are inherent to humans. The great
legacy of the document is the idea that individual rights come not from any
government, but from "their Creator." This ideal has led citizens to sometimes
reject a government or law that they feel violates their rights.

Legislation, Supreme Court cases, and amendments to the Constitution have all
been justified through interpretation of these documents. These government
actions, built on a foundation of constitutional principles, have expanded civil
rights to more people a little at a time.

\subsubsection{Civil Rights Expansion}

By the mid-1820s, the number of eligible voters grew, since most states had
dropped the requirement of property ownership in order to vote. Large groups of
people were still excluded though, like African Americans, who were enslaved in
the South and were not recognized as legal U.S. citizens. Women could not vote,
either. These groups also faced discrimination in public spaces and in
employment.

The Civil War was a bitter war, fought from 1861 to 1865 in the United States.
The Northern Union and Southern Confederate sides disagreed over issues such as
slavery and the rights of state government. In the five years following the war,
during what is known as Reconstruction, states ratified the 13th, 14th, and 15th
amendments. These three amendments changed the Constitution to establish former
slaves as U.S. citizens, giving them the same protections and rights as other
citizens. The amendments would lead to other legislation, amendments, and court
cases that targeted the civil rights of African Americans but influenced the
rights of other groups of people, as well.

\subsubsection{Incorporation}

The Civil Rights Movement and its associated landmark cases and legislation are
associated most with African Americans; however, other groups such as women and
Hispanic Americans benefited from the reforms, as well. Additionally, the
Supreme Court made other rulings in the 20th century that expanded civil rights
to other groups. A key feature binding many of these rulings together is they
lead to incorporation of the Bill of Rights.

In an 1833 case, the Supreme Court ruled that the Bill of Rights applied only to
the federal government, not the states. This began to change with various cases
in the 1900s, especially around the time of the Civil Rights Movement and the
Warren Court. Piece by piece, the Supreme Court ruled in a series of cases that
state governments also had to uphold most of the rights protected in the Bill of
Rights. They based their decisions on the "due process clause" of the 14th
Amendment that states in part, "Nor shall any State deprive any person of life,
liberty, or property, without due process of law." In other words, a state may
not place a person in jail, take away their rights, or take their property
without law enforcement following certain procedures to ensure fair conduct.
Thus, the process of incorporation further expanded civil rights to citizens of
all states.

\newpage
